%!TEX root = ../main.tex
\begin{center}
    \large
    \begin{singlespace}    
        \textbf{\chineseTitle{}} \\[0.5cm]
    \end{singlespace}

    \begin{singlespace}    
    學生：\studentCnName  \hspace{2.5cm}  指導教授：\advisorCnName \hspace{0.1cm} 博士 \\

    \ifdefined\advisorCnNameB % 如果有共同指導教授
        \hspace{9.6cm}  \advisorCnNameB ~ 博士 \\
    \fi
    \end{singlespace}

    \vspace{0.5cm}

    \begin{singlespace}
        國立陽明交通大學\\
        \NameofDepartmentInstituteCN\\[0.2cm]
    \end{singlespace}

    \textbf{摘\hspace{1cm}要} \\[0.5cm]

\end{center}

\normalsize 

基於深度學習的自動調變分類器（Automatic Modulation Classification, AMC）
在良好通道條件下可達到高分類準確率，但對於對抗式擾動（adversarial
perturbations）卻極為脆弱。以最佳化為基礎的攻擊方法，如
Carlini--Wagner（$L_2$）攻擊（CW）與彈性網路攻擊
（Elastic-Net Attack on Deep Neural Networks, EAD），
即使在原始波形大致仍可被解調的情況下，仍能顯著降低自適應小波網路
（Adaptive Wavelet Network, AWN）分類器於 RML2016.10a 資料集上的
準確率。本論文提出 \emph{Adaptive-K}，一種於接收端進行、
以頻域輸入轉換為基礎的防禦方法，可在不需重新訓練分類器的情況下
恢復其分類準確率。

Adaptive-K 在單次共用之快速傅立葉轉換（FFT）流程中，
整合三項機制：針對寬頻訊號的頻譜平坦度路由
（spectral-flatness routing）、可依訊噪比（SNR）自適應調整之
$K$ 值上限（K-cap）以平衡擾動抑制與訊號保留，
以及以每筆樣本之幅度轉折點（magnitude knee）估計所保留頻譜成分數量之方法。
本方法能保留主要承載調變資訊之頻譜成分，同時抑制由攻擊所引入之非主要頻譜雜訊。

我們針對 CW、EAD-L1、EAD-EN、FGSM 與 PGD 共五種攻擊，
於 RML2016.10a 資料集上以預訓練之 AWN 分類器評估 Adaptive-K 之效能，
並與五種依訊噪比逐一校正之傳統濾波器（Kalman、Wiener、Savitzky-Golay、
Gaussian、FIR）以及一項經驗性隨機平滑
（randomized smoothing）基準（$k = 20$ 份雜訊副本、$\sigma = 0.01$；
未提供形式化認證）進行比較。於 SNR $\geq 0$\,dB 之條件下，
CW 攻擊將加權平均準確率降至 35.4\%，EAD-L1 降至 6.6\%；
Adaptive-K 則將 CW 之準確率恢復至 71.3\%，EAD-L1 恢復至 70.8\%，
分別較最強之傳統濾波基準高出 6.2 與 8.5 個百分點。
對於 PGD 攻擊，其準確率由 32.6\% 提升至 43.4\%。
在乾淨訊號上之準確率下降不超過 0.5 個百分點，
且於本文之 GPU 實作中，每批 64 筆樣本之額外延遲低於 0.1\,ms。

以上結果顯示，自適應頻譜濾波為一種輕量且有效之防禦方法，
可於未知防禦資訊之對抗式攻擊下恢復 AMC 之分類效能。
本研究之評估並未主張可抵禦已知防禦資訊之自適應攻擊
（adaptive defense-aware attacks），此類攻擊之評估將為未來重要之研究方向。

\vspace{1cm}

% 中文摘要及關鍵詞 5-7 個
\textbf{關鍵字：}自動調變分類, 對抗式強健性, 頻譜防禦, FFT Top-$K$, RML2016.10a, 訊號處理
